% ------------------------------------------------------------------
% BIO 151 Live Lab — QR code for the student data-entry form.
% Paste the PREAMBLE block into the packet's preamble (Overleaf, pdfLaTeX
% or XeLaTeX both work) and the BODY block where the QR should appear.
% Change ONLY \LabFormURL if the deployment URL ever changes; the QR code,
% the printed link and the hyperlink all follow it automatically.
% ------------------------------------------------------------------

% ---------- PREAMBLE ----------
\usepackage{qrcode}
\usepackage{hyperref}
\usepackage{url}

% Verified student-form URL (GitHub Pages, HTTPS). No trailing text after the slash.
\newcommand{\LabFormURL}{https://justin83.github.io/biology-at-scc/labs/bio151-spectro/}

% Reusable box: QR on the left, instructions on the right.
\newcommand{\LabFormQR}{%
  \begin{center}
  \fbox{\begin{minipage}{0.92\linewidth}\vspace{4pt}
    \begin{minipage}[c]{1.25in}\centering
      \qrcode[height=1.1in,level=M]{\LabFormURL}
    \end{minipage}\hfill
    \begin{minipage}[c]{\dimexpr\linewidth-1.4in\relax}
      \textbf{Enter your readings on your phone.}\\
      Scan the code or go to \url{\LabFormURL}\\[3pt]
      \small 1.\ Pick your group number.\quad
      2.\ Type both partners' names once.\quad
      3.\ The first partner makes up a group PIN (4+ characters) and shares it; the other partner joins with the same PIN.\quad
      4.\ Type each absorbance as the display shows it and tap \textbf{Save}. You can come back and change any value.
    \end{minipage}\vspace{4pt}
  \end{minipage}}
  \end{center}}

% ---------- BODY ----------
% Recommended placement: at the top of the "Data Collection" / "Procedure Part B"
% page, immediately before the absorbance data table students would otherwise
% fill in by hand, so the phone form and the paper table sit side by side.
% A second, smaller copy near the title page is optional:
%   \qrcode[height=0.8in]{\LabFormURL}
\LabFormQR
